\documentclass[10pt,a4paper]{amsart}
\usepackage[margin=19mm]{geometry}
\usepackage{amsmath,amssymb,mathtools}
\usepackage[scr=rsfs,cal=euler]{mathalfa}
\usepackage{xfrac}
\usepackage[unicode,hidelinks,bookmarks=false]{hyperref}
\newcommand{\ie}{i.e.\ }
\newcommand{\cf}{cf.\ }
\newcommand{\resp}{respectively\ }
\newcommand{\Subsection}{\subsection}
\providecommand{\nobreakdash}{\nobreak\hbox{-}}
\setlength{\emergencystretch}{2em}
\multlinegap0pt
\usepackage{amssymb}
\newtheorem{lemma}{Lemma}
\usepackage{cleveref}
\crefname{lemma}{Lemma}{Lemmas}
\newcommand{\dd}{\mathop{}\!\mathrm{d}}
\title{Classification of solutions to the Liouville equation with a nonlinear Neumann boundary condition and applications to sharp inequalities}
\author{Xiaohan Cai and Abdolhakim SHOUMAN}
\date{Source: arXiv:2609.18762v1 (CC BY 4.0)}
\begin{document}
\maketitle

\noindent\emph{Source and licence.} Classification of solutions to the Liouville equation with a nonlinear Neumann boundary condition and applications to sharp inequalities. Version arXiv:2609.18762v1 (CC BY 4.0), distributed by arXiv under the Creative Commons Attribution 4.0 International licence, http://creativecommons.org/licenses/by/4.0/, as recorded in the arXiv licence metadata for this exact version and shown on the abstract page https://arxiv.org/abs/2609.18762v1. Full licence notice: \texttt{licenses/arxiv-2609.18762-CC-BY-4.0.txt}.\par\smallskip

\noindent\textit{Editorial scope.} Counted unit: the article's lemma ball (source0.tex lines 1598--1614). The article's complete original proof of it is delivered with it (source0.tex lines 1682--1949, one \texttt{proof} environment of 268 lines, which the article places away from the statement). No mathematics is written, repaired or re-derived: the statement and the proof are the article's own lines, in the article's own notation, and no retained line is substituted or re-flowed.  Retained uncounted as the source-local context the proof needs: lines 981--1009; lines 1616--1634.  Nothing was omitted from any retained span, so the package carries no omission record.  No retained line contains a citation, so the wrapper carries no bibliography and no bibliography entry.  No retained line contains a figure environment, an image-inclusion command or drawing code, so no image is delivered and the package declares no asset dependency.  Editorial formatting only, in addition: an AMS \texttt{amsart} preamble that restates the article's own declarations and macros used by the retained lines, namely the environments \texttt{lemma} on separate counters, together with the standard packages those lines need.  The retained environments print in this document as Lemma 1 in order of appearance, whereas the article prints the same items with the numbering of its own full document.  The complete native source of the article as distributed in the arXiv e-print for this exact version is bundled unchanged as \texttt{sources/210722/source0.tex}.
    \begin{align}
        \log\left(
        \frac{1}{\pi}\int_{\mathbb{B}^2}e^{2u}
        \right)
        \leq& \frac{1}{2\pi}\int_{\mathbb{B}^2}
        |\nabla u|^2
        +\frac{1}{\pi}
        \int_{\partial\mathbb{B}^2} u
        +\log\frac{eC_1}{8\pi},\quad 
        \forall u\in H^1(\mathbb{B}^2).
        \label{eq. weak Moser on ball}
        \\[3mm]
        \log\left(
        \frac{1}{\pi}\int_{\mathbb{B}^2}e^{2u}
        \right)
        \leq& \frac{1}{4\pi}\int_{\mathbb{B}^2}
        |\nabla u|^2+\log\frac{C_2}{\pi},\quad,\forall u\in H_0^1(\mathbb{B}^2).
        \label{eq. weak Moser for Dirichlet boundary value}
        \\[3mm]
        \log\left(
        \frac{1}{2\pi}
        \int_{\partial \mathbb{B}^2}e^{u}
        \right)
        \leq& \frac{1+\epsilon}{4\pi}\int_{\mathbb{B}^2}
        |\nabla u|^2
        +\frac{1}{2\pi}
        \int_{\partial\mathbb{B}^2} u
        +C_\epsilon,\quad \forall u\in H^1(\mathbb{B}^2),\qquad \forall \epsilon>0.\label{eq. weak LMO inequality}
    \end{align}

    \begin{lemma}\label{lem. concentration for measures}
        There exist two universal constants $\delta_0, d_0>0$ such that, for any probability measure $\mu$ with 
        \[
        \int_{\mathbb{B}^2}z\dd \mu=0,
        \]
        the following alternative holds
        \begin{enumerate}
            \item Either there holds
            \[
            \mu(B_{1-d_0})\geq \delta_0;
            \]
            \item or there exists $\Omega_1,\Omega_2\subset \overline{\mathbb{B}^2}$, such that $\operatorname{dist}(\Omega_1,\Omega_2)\geq d_0$ and
            \[
            \mu(\Omega_1)\geq \delta_0,
            \quad 
            \mu(\Omega_2)\geq \delta_0.
            \]
        \end{enumerate}
    \end{lemma}

    \begin{lemma}\label{lem. improved weak Moser on ball}
        For any $\epsilon>0$, there exists a constant $C_\epsilon>0$ such that, for any $u\in H^1(\mathbb{B}^2)$ with 
        \[
        \int_{\mathbb{B}^2}ze^{2u(z)}=0,
        \]
        there holds
        \[
        \log\left(
        \frac{1}{\pi}\int_{\mathbb{B}^2}e^{2u}
        \right)
        \leq \frac{1+\epsilon}{4\pi}\int_{\mathbb{B}^2}
        |\nabla u|^2
        +\frac{1}{\pi}
        \int_{\partial\mathbb{B}^2} u
        +C_\epsilon.
        \]
    \end{lemma}

    \begin{proof}[Proof of \cref{lem. improved weak Moser on ball}]
    {\bf Step 1:} We first assume that
    \[
    \int_{\mathbb{B}^2}u=0.
    \]
    For any $\epsilon>0$, let $V_\epsilon$ be the Neumann eigenspace corresponding to eigenvalues less than $\epsilon^{-1}$. Let $u_1\in L^{\infty}(\mathbb{B}^2)$ be the orthogonal projection of $u$ on $V_\epsilon$, and $u_2:=u-u_1\in H^1(\mathbb{B}^2)$. By definition and the Poincar\'e inequality, there holds
    \begin{align}
        ||u_1||_{L^{\infty}(\mathbb{B}^2)}
        \leq& C_\epsilon 
        ||u_1||_{L^2(\mathbb{B}^2)}
        \leq C_\epsilon ||\nabla u_1||_{L^2(\mathbb{B}^2)}
        \leq \epsilon\int_{\mathbb{B}^2}
        |\nabla u_1|^2
        +C_\epsilon,\label{eq. upper bound of u1}\\
        \int_{\mathbb{B}^2}|u_2|^2
        \leq& \epsilon 
        \int_{\mathbb{B}^2}|\nabla u_2|^2.
        \label{eq. L2 bound of u2}
    \end{align}

    
        By \cref{lem. concentration for measures}, the probability measure
        \[
        \dd\mu:=\frac{e^{2u(x)}}{\int_{\mathbb{B}^2}e^{2u}}\dd x
        \]
        satisfies the alternatives.

        {\bf Step 2:} For the first alternative, i.e.
        \[
        \int_{B_{1-d_0}}e^{2u}
        \geq \delta_0\int_{\mathbb{B}^2}e^{2u}.
        \]
        Let $\chi\in C_0^{\infty}(\mathbb{B}^2)$ be a cut-off function with $\chi\equiv 1$ in $B_{1-d_0}$. Then we have
        \begin{equation}\label{eq. area upper bound in first alternative}
            \delta_0 \int_{\mathbb{B}^2}e^{2u}
        \leq \int_{B_{1-d_0}}e^{2u}
        \leq e^{2||u_1||
        _{L^\infty(\mathbb{B}^2)}}
        \int_{\mathbb{B}^2}e^{2\chi u_2}.
        \end{equation}
        By Cauchy-Schwarz inequality and \cref{eq. L2 bound of u2}, there holds, for any $\eta>0$,
        \begin{equation}
        \label{eq. H01 norm upper bound}
            \int_{\mathbb{B}^2}|\nabla(\chi u_2)|^2
        \leq (1+\eta)
        \int_{\mathbb{B}^2}|\nabla u_2|^2
        +C_\eta\int_{\mathbb{B}^2}|u_2|^2
        \leq (1+\eta+C_\eta\epsilon)
        \int_{\mathbb{B}^2}|\nabla u_2|^2
        \end{equation}
        Applying \cref{eq. weak Moser for Dirichlet boundary value} to $\chi u_2$ and using \cref{eq. H01 norm upper bound}, we obtain
        \[
        \log\left(
        \frac{1}{\pi}
        \int_{\mathbb{B}^2}e^{2\chi u_2}
        \right)
        \leq \frac{1+\eta+C_\eta\epsilon}{4\pi}
        \int_{\mathbb{B}^2}|\nabla u_2|^2
        +C.
        \]
        Inserting this into \cref{eq. area upper bound in first alternative}, taking logarithm and using \cref{eq. upper bound of u1}, we get 
        \begin{align*}
            \log\left(
        \frac{1}{\pi}
        \int_{\mathbb{B}^2}e^{2u}
        \right)
        \leq& 2\epsilon\int_{\mathbb{B}^2}
        |\nabla u_1|^2
        +C_\epsilon
        +\frac{1+\eta+C_\eta\epsilon}{4\pi}
        \int_{\mathbb{B}^2}|\nabla u_2|^2
        +C\\
        \leq &\frac{1+\eta+C_\eta\epsilon}{4\pi}
        \int_{\mathbb{B}^2}|\nabla u|^2
        +C_\epsilon.
        \end{align*}
        First fixing $\eta\ll 1$ and then choosing $\epsilon\ll 1$, after renaming the parameters, we get
        \[
        \log\left(
        \frac{1}{\pi}
        \int_{\mathbb{B}^2}e^{2u}
        \right)
        \leq \frac{1+\epsilon}{4\pi}
        \int_{\mathbb{B}^2}|\nabla u|^2
        +C_\epsilon.
        \]    

        {\bf Step 3:} For the second alternative, i.e. there exists $\Omega_1,\Omega_2\subset \overline{\mathbb{B}^2}$, such that $\operatorname{dist}(\Omega_1,\Omega_2)\geq d_0$ and
        \[
        \int_{\Omega_1}e^{2u}
        \geq \delta_0
        \int_{\mathbb{B}^2}e^{2u},\quad
        \int_{\Omega_2}e^{2u}
        \geq \delta_0
        \int_{\mathbb{B}^2}e^{2u}.
        \]
        Let $\chi_1,\chi_2\in C^\infty(\mathbb{B}^2)$ be two functions with disjoint support and $\chi_i\equiv 1$ on $\Omega_i$. Then we have
        \begin{equation}\label{eq. area upper bound in second alternative}
            \delta_0 \int_{\mathbb{B}^2}e^{2u}
        \leq \int_{\Omega_i}e^{2u}
        \leq e^{2||u_1||
        _{L^\infty(\mathbb{B}^2)}}
        \int_{\mathbb{B}^2}
        e^{2\chi_i u_2}.
        \end{equation}
        Note that, by \cref{eq. L2 bound of u2},
        \begin{align}
            \int_{\operatorname{supp}\chi_i}
            |\nabla(\chi_i u_2)|^2
        \leq& (1+\eta)
        \int_{\operatorname{supp}\chi_i}
        |\nabla u_2|^2
        +C_\eta\int_{\mathbb{B}^2}|u_2|^2\notag
        \\
        \leq& (1+\eta)
        \int_{\operatorname{supp}\chi_i}|\nabla u_2|^2
        +C_\eta\epsilon
        \int_{\mathbb{B}^2}|\nabla u_2|^2.
        \label{eq. H01 norm upper bound in second}
        \end{align}
        By Sobolev-trace inequality, \cref{eq. L2 bound of u2}, \cref{eq. H01 norm upper bound in second} and Young's inequality, there holds
        \begin{align}
        \label{eq. upper bound of trace}
            \int_{\partial\mathbb{B}^2}\chi_i u_2
            \leq \sqrt{2\pi}
            ||\chi_2u_2||
            _{L^2(\partial \mathbb{B}^2)}
            \leq \sqrt{2\pi}
            ||\chi_2u_2||
            _{H^{1}(\mathbb{B}^2)}
            \leq \epsilon \int_{\mathbb{B}^2}
            |\nabla u_2|^2+C_\epsilon.
        \end{align}
        

        Applying \cref{eq. weak Moser on ball} to $\chi_i u_2$ and inserting \cref{eq. H01 norm upper bound in second} and \cref{eq. upper bound of trace} into it, we obtain
        \[
        \log\left(
        \frac{1}{\pi}
        \int_{\mathbb{B}^2}e^{2\chi_i u_2}
        \right)
        \leq 
        \frac{1+\eta}{2\pi}
        \int_{\operatorname{supp}\chi_i}|\nabla u_2|^2
        +\left(\frac{C_\eta\epsilon}{2\pi}
        +\frac{\epsilon}{\pi}
        \right)
        \int_{\mathbb{B}^2}|\nabla u_2|^2
            +C_\eta+C_\epsilon.
        \]
        Inserting this into \cref{eq. area upper bound in second alternative}, taking logarithm and using \cref{eq. upper bound of u1}, we get 
        \[
        \log\left(
        \frac{1}{\pi}
        \int_{\mathbb{B}^2}e^{2u}
        \right)
        \leq
        2\epsilon\int_{\mathbb{B}^2}
        |\nabla u_1|^2
        +\frac{1+\eta}{2\pi}
        \int_{\operatorname{supp}\chi_i}|\nabla u_2|^2
        +\left(\frac{C_\eta\epsilon}{2\pi}
        +\frac{\epsilon}{\pi}
        \right)
        \int_{\mathbb{B}^2}|\nabla u_2|^2
        +C_\eta+C_\epsilon.
        \]
        By adding the inequalities corresponding to $i=1$ and $i=2$, we derive
        \begin{align*}
            2\log\left(
        \frac{1}{\pi}
        \int_{\mathbb{B}^2}e^{2u}
        \right)
        &\leq
        4\epsilon\int_{\mathbb{B}^2}
        |\nabla u_1|^2
        +\frac{1+\eta}{2\pi}
        \int_{\mathbb{B}^2}|\nabla u_2|^2
        +\left(\frac{C_\eta\epsilon}{\pi}
        +\frac{2\epsilon}{\pi}
        \right)
        \int_{\mathbb{B}^2}|\nabla u_2|^2
        +C_\eta+C_\epsilon\\[3mm]
        &\leq 
        \left(
        \frac{1+\eta}{2\pi}
        +\frac{(C_\eta+2)\epsilon}{\pi}
        \right)
        \int_{\mathbb{B}^2}|\nabla u|^2
        +C_\eta+C_\epsilon.
        \end{align*}
        First fixing $\eta\ll 1$ and then choosing $\epsilon\ll 1$, after renaming the parameters, we get
        \[
        \log\left(
        \frac{1}{\pi}
        \int_{\mathbb{B}^2}e^{2u}
        \right)
        \leq \frac{1+\epsilon}{4\pi}
        \int_{\mathbb{B}^2}|\nabla u|^2
        +C_\epsilon.
        \] 

        {\bf Step 4:} We have proved that, for any $\epsilon>0$, there exists $C_\epsilon>0$, such that if $\int_{\mathbb{B}^2}ze^{2u(z)}=0$ and
        $\int_{\mathbb{B}^2}u=0$, then  
        \[
        \log\left(
        \frac{1}{\pi}
        \int_{\mathbb{B}^2}e^{2u}
        \right)
        \leq \frac{1+\epsilon}{4\pi}
        \int_{\mathbb{B}^2}|\nabla u|^2
        +C_\epsilon.
        \]
        By translation, this is equivalent to
        \begin{equation}
        \label{eq. improved weak Moser on ball with interior mean}
            \log\left(
        \frac{1}{\pi}
        \int_{\mathbb{B}^2}e^{2u}
        \right)
        \leq \frac{1+\epsilon}{4\pi}
        \int_{\mathbb{B}^2}|\nabla u|^2
        +\frac{2}{\pi}\int_{\mathbb{B}^2}u
        +C_\epsilon,\quad 
        \forall u\in H^1(\mathbb{B}^2)
        \text{ with } 
        \int_{\mathbb{B}^2}ze^{2u(z)}=0.
        \end{equation}
        
        Now we claim that
        \begin{equation}
        \label{eq. difference between interior mean and boundary mean}
            \left|
        \frac{2}{\pi}\int_{\mathbb{B}^2}u
        -\frac{1}{\pi}
        \int_{\partial\mathbb{B}^2}u
        \right|
        \leq C
        \left(\int_{\mathbb{B}^2}|\nabla u|^2
        \right)^{\frac{1}{2}},\quad 
        \forall u\in H^1(\mathbb{B}^2).
        \end{equation}
        Otherwise, there exists a sequence $u_j\in H^1(\mathbb{B}^2)$ such that
        \begin{equation}
        \label{eq. blow up condition}
            \int_{\mathbb{B}^2}|\nabla u_j|^2=1,\quad 
        \left|
        \frac{2}{\pi}\int_{\mathbb{B}^2}u_j
        -\frac{1}{\pi}
        \int_{\partial\mathbb{B}^2}u_j
        \right|
        \geq j.
        \end{equation}
        After a translation, we assume that $\int_{\mathbb{B}^2}u_j=0$. Then by \cref{eq. blow up condition}, Cauchy-Schwarz inequality, Sobolev-trace inequality and Poincar\'e inequality, we have
        \[
        j
        \leq \frac{1}{\pi}\int_{\partial \mathbb{B}^2}
        |u_j|
        \leq 2\left(
        \frac{1}{2\pi}
        \int_{\partial\mathbb{B}^2}|u_j|^2
        \right)^{\frac{1}{2}}
        \leq C||u_j||_{H^1(\mathbb{B}^2)}
        \leq C||\nabla u_j||_{L^2(\mathbb{B}^2)}
        =C.
        \]
        This is a contradiction. Combining \cref{eq. improved weak Moser on ball with interior mean} with \cref{eq. difference between interior mean and boundary mean} completes the proof.
    \end{proof}

\end{document}
